\chapter[Referencial Teórico]{Referencial Teórico}

Neste capítulo serão apresentados os conceitos chaves, teorias, modelos que serão necessários para compreensão completa da pesquisa.


\section{Engenharia de software Experimental}

A engenharia de software tornou-se hoje uma área muito influencial ao comportamento humano, isso prova-se da forma como existem inúmeras organizações com milhares de variações de fluxo de desenvolvimento. Nesse cenário, a mera replicação de um modelo de processo bem-sucedido não assegura, de forma isolada, a obtenção de resultados equivalentes em contextos distintos. Visando mitigar variabilidade em projetos de sofware, várias decisões devem ser tomadas tais como, seleção de ferramentas, formalização de processos e padronização de métodos e práticas. Portanto, o formato não determinístico da constução de um software torna-se próxima de uma realidade mensurável. Sendo assim por meio da engenharia de softwarea experimental podemos aprofundar e trazer padrões necessários para definições de um objetivo válido, ela traz de forma sistemática, controlada e quantificável de avaliar as atividades da área, mesmo que sejam baseadas em comportamentos humanos. \cite{wohlin2024experimentation}


\subsection{Revisão Estruturada da Literatura}

Para fundamentar pesquisas baseadas em evidências empíricas, a Revisão Estruturada da Literatura é um método utilizado para identificar, avaliar e interpretar pesquisas disponíveis e relevantes sobre uma questão específica ou fenômeno de interesse. Diferentemente de uma revisão bibliográfica ad-hoc, a mesma segue protocolos bem estabelecidos, reduzindo vieses durante o processo de análise. \cite{wohlin2024experimentation}.

O protocolo da Revisão consiste no atendimento dos seguintes:

\begin{itemize}
    \item Planejamento da revisão 
    \begin{itemize}
        \item Definição da pergunta principal de pesquisa e perguntas secundárias;
        \item Montagem da String de busca utilizando o protocolo PICO;
        \item Definição de critérios de inclusão e exclusão;
        \item Definição do formulário de extração de dados;
    \end{itemize}
    \item Condução
    \begin{itemize}
        \item Seleção dos estudos com base nos critérios de inclusão e exclusão;
        \item Avaliação da qualidade dos estudos selecionados;
        \item Extração dos dados relevantes dos estudos selecionados;
    \end{itemize}
    \item Síntese dos resultados
    \begin{itemize}
        \item Análise qualitativa e/ou quantitativa dos dados extraídos;
        \item Interpretação dos resultados em relação à questão de pesquisa;
    \end{itemize}
\end{itemize}


\subsubsection{Protocolo PICO}

Na etapa de planejamento da revisão estruturada, a string de busca não sendo definida dentro do contexto da pergunta de pesquisa, pode trazer dificuldades ao pesquisador em encontrar os artigos relevantes ao problema investigado. Dessa forma o protocolo PICO auxilia na construção de uma string de busca mais fiel ao escopo da pesquisa. O acrônimo PICO é composto por quatro elementos: População, Intervenção, Comparação e Resultados (\textit{Outcomes}). \cite{pai2004systematic}

\subsection{Estudo de caso}

O estudo de caso é uma investigação empírica de um fenômeno contemporâneo, ou vários pequenas ocorrências dele dentro de um contexto real. Esse tipo de estudo de estudo é adequado para avaliar métodos e ferramentas de engenharia de software, pois permite ao pesquisadaor observar o fenômeno em tempo real. Devido a esse contexto, seu planejamento, coleta de análise de dados seguem uma abordagem flexível, oque significa que essas etapas acontecem de forma iterativa e são evoluidas com a maturidade da pesquisa.

Entretanto sua flexibilidade não deixa de ser método de pesquisa, portanto, exige um protocolo a ser seguido, desde sua definição, planejamento, execução e análise dos dados. O protocolo é um documento que descreve o plano de estudo de caso, incluindo os objetivos, questões de pesquisa, hipóteses, variáveis, métodos de coleta e análise de dados, cronograma e recursos necessários. \cite{yin2001estudo}

\begin{itemize}
    \item \textbf{1. Visão Geral do Projeto do Estudo de Caso}
    \begin{itemize}
        \item Identificação do Projeto
        \item Contexto
        \item Objetivo Principal
        \item Questões e Proposições do Estudo
        \item Leituras e Normativos Relevantes
    \end{itemize}
    
    \item \textbf{2. Procedimentos de Campo}
    \begin{itemize}
        \item Estratégia de Acesso aos Dados
        \item Garantias de Segurança e Ética
        \item Preparação do Ambiente de Análise
        \item Prazos e Cronograma
    \end{itemize}
    
    \item \textbf{3. Questões do Estudo de Caso}
    \begin{itemize}
        \item Questões sobre o Contrato
        \item Questões sobre o Produto 
        \item Planilhas e Matrizes de Coleta de Dados
    \end{itemize}
    
    \item \textbf{4. Guia para o Relatório do Estudo de Caso}
    \begin{itemize}
        \item Público-alvo
        \item Formato de Apresentação
        \item Tratamento das Evidências
        \item Composição dos Resultados e Generalização Analítica
    \end{itemize}
\end{itemize}

\subsection{Entrevistas}

No Estudo de caso há a possibiliade na coleta de dados ser  baseada em entrevistas, o pesquisador faz uma série de perguntas a um conjunto de pessoas sobre as áreas de interesse no estudo de caso. Na maioria dos casos, uma entrevista é realizada com cada pessoa individualmente, mas é possível realizar entrevistas em grupo, por exemplo, grupos focais. O diálogo entre o pesquisador e entrevistado é guiado por um conjunto de perguntas de entrevista. \cite{wohlin2024experimentation}

As perguntas da entrevista são baseadas nas perguntas da pesquisa ou que seja no enfoque da pesquisa principal. As perguntas podem ser abertas, ou seja, permitindo e convidando a uma ampla gama de respostas e questões do entrevistado, ou fechadas, ou seja, oferecendo um conjunto limitado de respostas alternativas.



\section{Modelo de Contratação de Serviços de Desenvolvimento de Software}

No Brasil, a terceirização de serviços de TI foi ganhando notoriedade com denominação do Sistema de Administração dos Recursos de Informação e Informática (SISP) por meio do decreto 1.048 de 21 de Janeiro de 1994 revogado pelo Decreto nº 7.579 de 11 de Outubro de 2011 \cite{brasil2011decreto7579}, essa desde então é a espínha dorsal da governança de TI no poder público federal. É por meio desta que teve a primeira instrução normativa para contratação de desenvolvimento de software IN 04 da SLTI/MPOG, que teve várias alterações devido a complexidade de software, a primeira publicação foi no ano de 2008 e teve aperfeiçoamento até 2014.

Posteriormente devido a mudança estrutural da SLTI para a SGD determinada pelos decretos Decreto nº 8.578 \cite{brasil2015decreto8578}, Decreto nº 9.679 e 9.745 \cite{brasil2019decreto9679, brasil2019decreto9745} , onde o mesmo ficaria sob o Ministério da Economia e que posteriormente a publicação do Decreto n 11.345/2023 reposiciona a SGD sob o Ministério da Gestão e da Inovação em Serviços Públicos. Nessa última Houve a alteração da IN 04/2014 para a IN 1/2019. 

Por fim, a vinda da nova lei de licitação de contratos em 2021, teve a alteração da IN 1/2019 para a IN 94/2022 \cite{brasil2022in94}, que essa é usado como geral e base para a publicação do modelo de contratação de serviços de desenvolvimento, Manutenção e Sustentação de Software, por meio da portaria nº 750/2023 da SGD \cite{brasil2023portaria750}.

No contexto do Governo do Distrito Federal, a normatização das contratações de Tecnologia da Informação alinhou-se ao rigor do ecossistema federal por meio do Decreto Distrital nº 44.330/2023 \cite{distritofederal2023decreto44330}, que instituiu a recepção subsidiária dos atos normativos da União. Consequentemente, ao absorver a Instrução Normativa SGD/ME nº 94/2022, o arcabouço jurídico distrital tornou de observância obrigatória as diretrizes da Portaria SGD/MGI nº 750/2023 para os serviços de desenvolvimento, manutenção e sustentação de software.

\section{Modelos de Qualidade de Software}

A avaliação e garantia da qualidade de um software exigem a detecção e medição de diversos parâmetros, características e sub-características. A pesquisa sobre a qualidade de software é tão antiga quanto a própria construção de software, e o uso de modelos estabelecidos é o meio mais aceito para apoiar a gestão da qualidade do produto \cite{miguel2014review}. Um modelo de qualidade pode ser definido como um conjunto de características e as relações entre elas, fornecendo a base para a especificação de requisitos de qualidade e sua respectiva avaliação \cite{miguel2014review}.

Historicamente, a literatura consolidou diversos modelos hierárquicos, como o Modelo de McCall (1977), o Modelo de Boehm (1978), FURPS (1987) e Dromey (1995) \cite{curcio2016}. O modelo FURPS, por exemplo, dividiu as características em requisitos funcionais e não-funcionais (Usabilidade, Confiabilidade, Desempenho e Suportabilidade) \cite{miguel2014review}. Posteriormente, a \textit{International Organization for Standardization} (ISO) padronizou o conceito de qualidade de produto de software com a publicação da norma ISO/IEC 9126, que definiu seis características principais: funcionalidade, confiabilidade, usabilidade, eficiência, manutenibilidade e portabilidade \cite{curcio2016}.

Devido à evolução das engenharias de software, a ISO/IEC 9126 foi atualizada e substituída pela família de normas ISO/IEC 25000, conhecida como SQuaRE (\textit{Software Product Quality Requirements and Evaluation}) \cite{curcio2016}. Especificamente, a norma ISO/IEC 25010 redefiniu o modelo de qualidade em oito características fundamentais: Adequação Funcional, Confiabilidade, Eficiência de Desempenho, Operabilidade (Usabilidade), Segurança, Compatibilidade, Manutenibilidade e Transferibilidade \cite{miguel2014review}. 

\subsection{Modelo Q-Rapids}

Na busca por operacionalizar e automatizar a medição de qualidade e apoiar a tomada de decisão ágil, destaca-se o modelo Q-Rapids. Trata-se de um modelo de qualidade e uma ferramenta de \textit{software analytics} desenvolvida sob o consórcio do programa de pesquisa e inovação \textit{Horizon 2020} da União Europeia \cite{qrapids2019} [10]. 

O Q-Rapids propõe a avaliação contínua da qualidade do software por meio da coleta de dados brutos provenientes de ferramentas de uso contínuo no ciclo de desenvolvimento, como \textit{SonarQube}, \textit{Jira} e \textit{Jenkins} \cite{qrapids2019}. O modelo é estruturado de forma hierárquica em Indicadores Estratégicos e Fatores de Processo e Produto \cite{qrapids2019}.

Dentre os Indicadores Estratégicos propostos pelo modelo, destacam-se a "Qualidade do Produto" (\textit{Product Quality}) e o "Bloqueio" (\textit{Blocking}) \cite{qrapids2019}. O modelo Q-Rapids estabelece que a "Qualidade do Código" é um fator de produto aferido a partir de métricas como Complexidade Ciclomática, Comentários e Duplicação, extraídas diretamente do \textit{SonarQube} \cite{qrapids2019}. Além disso, introduz o conceito prático de "Código Bloqueante" (\textit{Blocking code}), que mede a Dívida Técnica do código de software em termos de violações de regras de qualidade, gerando ações imediatas, como a exigência de refatoração para violações críticas ou bloqueantes \cite{qrapids2019}.

\subsection{Modelo SQALE}

O objetivo principal do modelo SQALE é fornecer suporte à avaliação e gestão da Dívida Técnica, oferecendo métricas objetivas, precisas e reprodutíveis que refletem de forma transparente o estado estrutural do código \cite{letouzey2010}.

A arquitetura do método SQALE atua sobre duas hierarquias: a estrutura do código-fonte e o Modelo de Qualidade. O Modelo de Qualidade do SQALE organiza os requisitos do software de forma hierárquica (características, sub-características e regras de código), alinhando-se com a cronologia das atividades do ciclo de vida do desenvolvimento \cite{letouzey2010}. As características avaliadas dialogam diretamente com padrões internacionais, incluindo Testabilidade, Confiabilidade, Capacidade de Mudança (\textit{Changeability}) e Manutenibilidade.

O grande diferencial do SQALE reside em seu Modelo de Análise, que soluciona o problema de como agregar métricas de software de forma matematicamente válida. O modelo baseia-se no conceito de "Índice de Remediação". Avaliar o código-fonte, sob a ótica do SQALE, consiste em medir a distância que separa o código atual da sua meta ideal de qualidade \cite{letouzey2010}. Dessa forma, para cada violação de regra ou não-conformidade detectada no código, o modelo estima o esforço (geralmente quantificado em tempo, como minutos ou horas) necessário para corrigir o defeito e adequá-lo aos padrões \cite{letouzey2010}.

Como o índice de remediação representa estritamente um esforço de trabalho, a consolidação das métricas ocorre por meio de uma simples adição matemática através da árvore de componentes do software (ex: soma-se o esforço dos métodos para obter o da classe, das classes para o arquivo, e dos arquivos para o projeto completo) \cite{letouzey2010}. Essa propriedade garante que o método respeite rigorosamente a teoria da medição, evitando as distorções graves e perda de sensibilidade causadas pelo uso de médias ou percentuais em análises de qualidade tradicionais \cite{letouzey2010}.

\subsubsection{Dívida Técnica (Technical Debt)}

O termo Dívida Técnica (\textit{Technical Debt} - TD) foi originalmente cunhado por Ward Cunningham utilizando uma metáfora financeira para explicar os problemas associados à qualidade do código-fonte postergada \cite{silva2018}. Atualmente, a Dívida Técnica é definida de forma consolidada como um conjunto de construções de \textit{design} ou implementação que são convenientes a curto prazo (permitindo entregas rápidas), mas que configuram um contexto técnico que tornará mudanças futuras muito mais custosas ou até impossíveis \cite{silva2018}. Em suma, refere-se a atalhos e soluções alternativas que afetam principalmente a qualidade interna do software, em especial a sua manutenibilidade \cite{silva2018}.

Assim como no mercado financeiro, a Dívida Técnica possui um "principal", que é o esforço necessário para resolver a diferença entre o estado atual do código e a qualidade ideal, e os "juros", que refletem o esforço extra e o custo adicional cobrado da equipe sempre que ela precisar manter ou evoluir um software altamente acoplado ou complexo \cite{silva2018}. Se não gerenciado, o acúmulo excessivo dessa dívida pode levar o projeto a uma "falência técnica", estágio onde o sistema não pode mais evoluir de forma viável \cite{silva2018}.

Martin Fowler propôs uma taxonomia dividindo a TD em quatro quadrantes, focada na intencionalidade: pode ser deliberada (quando a equipe assume o risco por pressão de prazo) ou inadvertida (quando a equipe falha por falta de conhecimento); e pode ser imprudente (sem planos de pagar) ou prudente (estratégica e monitorada) \cite{silva2018}. Alves et al. ampliaram a classificação técnica, dividindo a dívida em tipos como: dívida de código (\textit{code debt}), dívida de testes (\textit{test debt}), dívida de \textit{design} e dívida de arquitetura \cite{silva2018}. 


